\documentclass[11pt]{article}

\usepackage[margin=1.05in]{geometry}
\usepackage{amsmath,amssymb,amsthm,mathtools}
\usepackage{booktabs}
\usepackage{enumitem}
\usepackage{microtype}
\usepackage{xcolor}
\usepackage[hidelinks]{hyperref}
\usepackage{fancyhdr}

\definecolor{deepblue}{HTML}{173B57}
\definecolor{warmgray}{HTML}{5E646A}
\hypersetup{
  colorlinks=true,
  linkcolor=deepblue,
  citecolor=deepblue,
  urlcolor=deepblue,
  pdftitle={The exact quantum value of a cyclic Bell operator},
  pdfauthor={Alec Kriebel}
}

\pagestyle{fancy}
\fancyhf{}
\fancyhead[L]{\small\textsc{Exact cyclic Bell bound}}
\fancyhead[R]{\small\textsc{A. Kriebel}}
\fancyfoot[C]{\thepage}
\setlength{\headheight}{14pt}

\newtheorem{theorem}{Theorem}[section]
\newtheorem{lemma}[theorem]{Lemma}
\newtheorem{proposition}[theorem]{Proposition}
\newtheorem{corollary}[theorem]{Corollary}
\theoremstyle{remark}
\newtheorem{remark}[theorem]{Remark}

\newcommand{\Id}{\mathbf 1}
\newcommand{\cI}{\mathcal I}
\newcommand{\cH}{\mathcal H}
\newcommand{\Reop}{\operatorname{Re}}
\newcommand{\Tr}{\operatorname{Tr}}
\newcommand{\spec}{\operatorname{spec}}
\newcommand{\dd}{\,\mathrm d}
\newcommand{\e}{\mathrm e}
\newcommand{\bbC}{\mathbb C}
\newcommand{\bbZ}{\mathbb Z}
\newcommand{\absop}[1]{\lvert #1\rvert}

\title{\vspace{-1.4cm}\textbf{The exact quantum value of a cyclic Bell operator}\\[0.45em]
\large A sharp polar-decomposition certificate for all dimensions}
\author{Alec Kriebel\\
\small with heavy assistance from ChatGPT 5.6 Sol}
\date{26 July 2026}

\begin{document}
\maketitle

\begin{abstract}
Perito, D'Avino, Jung, Mironowicz, Acin, and Augusiak recently
introduced a cyclic family of Bell operators and conjectured that its
Tsirelson bound is
\(
2\csc(\pi/(2d))
\)
for every integer \(d\geq2\).  We prove the conjecture.  In fact, the
operator upper bound holds after relaxing every observable to an arbitrary
unitary; the order-\(d\) relations are not needed.  The proof combines an
exact positive-factor identity from polar decomposition with the sharp scalar
formula
\[
  \max_{\lvert z\rvert=1}\sum_{y=0}^{d-1}\lvert1+\e^{2\pi i y/d}z\rvert
  =2\csc\!\left(\frac{\pi}{2d}\right).
\]
We characterize equality in the scalar formula, give an explicit
functional-calculus certificate for the operator inequality, and include a
self-contained verification that the known \(d\)-dimensional strategy is
admissible and saturates the bound.  This determines the quantum value only;
it does not establish uniqueness, self-testing, or all-dimensional
randomness certification.
\end{abstract}

\medskip
\noindent\textbf{Status.}
This is an unrefereed, AI-assisted research note.  The analytic argument is
accompanied by symbolic and numerical checks, but has not received independent
expert review.

\section{Introduction}

Let \(d\geq2\) and \(\omega=\e^{2\pi i/d}\).  For unitary observables
\(A_0,A_1\) for Alice and \(B_0,\ldots,B_{d-1}\) for Bob, Perito et
al.~\cite{PeritoEtAl2026} consider
\begin{equation}
  \cI_d
  =
  \frac12\sum_{x=0}^{1}\sum_{y=0}^{d-1}
       \omega^{xy} A_x\otimes B_y+\mathrm{h.c.}
  =
  \sum_{y=0}^{d-1}
       \Reop\!\left((A_0+\omega^yA_1)\otimes B_y\right).
  \label{eq:bell}
\end{equation}
In the Bell scenario the observables obey
\(
A_x^d=B_y^d=\Id
\).
The originating paper constructs a \(d\)-dimensional strategy with value
\begin{equation}
  \lambda_d=2\csc\!\left(\frac{\pi}{2d}\right),
  \label{eq:lambda}
\end{equation}
and conjectures that this is optimal.  It proves the weaker general upper
bound \(d\sqrt2\), and reports numerical tightness of
\eqref{eq:lambda} through \(d=6\).

Our contribution is the missing analytic upper bound.  The essential point is
that the Bob operators can be optimized term by term by polar decomposition
without destroying the coupling between \(A_0\) and \(A_1\).  That coupling is
then entirely encoded by the relative unitary \(A_0^\dagger A_1\), reducing
the problem to a scalar trigonometric extremum on the unit circle.

\begin{theorem}[Sharp arbitrary-unitary bound]
\label{thm:main}
Let \(d\geq2\).  On arbitrary complex Hilbert spaces, let
\(A_0,A_1,B_0,\ldots,B_{d-1}\) be unitaries, with the Alice and Bob
operators acting on separate tensor factors.  Then
\begin{equation}
  \cI_d\ \leq\ 2\csc\!\left(\frac{\pi}{2d}\right)\Id
  \label{eq:operator-bound}
\end{equation}
as an operator inequality.  Consequently, when the order-\(d\) relations are
imposed,
\[
  \boxed{\displaystyle
  \beta_{\mathrm Q}(\cI_d)
  =2\csc\!\left(\frac{\pi}{2d}\right)}
  \qquad(d\geq2).
\]
\end{theorem}

The upper bound therefore holds on a strict relaxation of the original
optimization problem.  It is nevertheless tight because the attaining
strategy of Ref.~\cite{PeritoEtAl2026} satisfies the order-\(d\) constraints.

\section{A polar-decomposition identity}

For a bounded operator \(C\), write
\(
\absop{C}=(C^\dagger C)^{1/2}
\)
and
\(
\absop{C^\dagger}=(CC^\dagger)^{1/2}
\).

\begin{lemma}[Polar positive-factor identity]
\label{lem:polar}
Let \(C\) be a bounded operator on \(\cH_A\), let
\(C=V\absop C\) be its polar decomposition, and let \(B\) be unitary on
\(\cH_B\).  Define
\[
  P_{C,B}
  =
  \absop{C^\dagger}^{1/2}\otimes\Id
  -
  V\absop C^{1/2}\otimes B .
\]
Then
\begin{equation}
  \frac{\absop C+\absop{C^\dagger}}2\otimes\Id
  -
  \Reop(C\otimes B)
  =
  \frac12 P_{C,B}^\dagger P_{C,B}.
  \label{eq:polar-sos}
\end{equation}
\end{lemma}

\begin{proof}
The polar factor \(V\) may be a partial isometry.  Its initial projection is
the support projection of \(\absop C\), so
\[
  \absop C^{1/2}V^\dagger V\absop C^{1/2}=\absop C.
\]
Polar functional calculus also gives
\[
  \absop{C^\dagger}^{1/2}V
  =
  V\absop C^{1/2},
  \qquad
  \absop{C^\dagger}^{1/2}V\absop C^{1/2}=C.
\]
Expanding \(P_{C,B}^\dagger P_{C,B}\) and using \(B^\dagger B=\Id\)
therefore yields
\[
  \absop C+\absop{C^\dagger}
  -C\otimes B-C^\dagger\otimes B^\dagger,
\]
with identity operators on the omitted tensor factors.  Dividing by two
proves \eqref{eq:polar-sos}.
\end{proof}

\begin{remark}
No invertibility or finite-dimensional unitary extension of \(V\) is needed.
This is useful because some of the operators \(A_0+\omega^yA_1\) may be
singular for general strategies.
\end{remark}

\section{The scalar roots-of-unity extremum}

\begin{lemma}[Sharp scalar bound and equality set]
\label{lem:scalar}
For every \(d\geq2\) and every \(z\in\bbC\) with \(\lvert z\rvert=1\),
\begin{equation}
  \sum_{y=0}^{d-1}\lvert1+\omega^y z\rvert
  \leq
  2\csc\!\left(\frac{\pi}{2d}\right).
  \label{eq:scalar}
\end{equation}
Equality holds if and only if
\begin{equation}
  z^d=(-1)^{d-1}.
  \label{eq:equality-set}
\end{equation}
\end{lemma}

\begin{proof}
Write \(z=\e^{2is}\), put
\(
h=\pi/d
\)
and
\(
a=h/2=\pi/(2d)
\).
Then the left side of \eqref{eq:scalar} is \(2G_d(s)\), where
\[
  G_d(s)=\sum_{y=0}^{d-1}\lvert\cos(s+yh)\rvert.
\]
The function \(G_d\) has period \(h\).

First suppose \(d=2m+1\) is odd.  We may take \(-a\leq s\leq a\)
and reindex the summands by \(k=-m,\ldots,m\).  Every angle
\(s+kh\) lies in \([-\pi/2,\pi/2]\), hence
\[
  G_d(s)
  =
  \sum_{k=-m}^{m}\cos(s+kh)
  =
  \frac{\sin(dh/2)}{\sin(h/2)}\cos s
  =
  \csc(a)\cos s.
\]
The maximum is \(\csc(a)\), attained exactly when
\(s\equiv0\pmod h\).  Equivalently, \(z^d=1=(-1)^{d-1}\).

Now suppose \(d=2m\) is even.  Take \(0\leq s\leq h\).  The first
\(m\) cosines are nonnegative and the last \(m\) are nonpositive, so
the finite cosine-sum formula gives
\begin{align*}
  G_d(s)
  &=
  \sum_{y=0}^{m-1}\cos(s+yh)
  -
  \sum_{y=m}^{2m-1}\cos(s+yh)\\
  &=
  \csc(a)\cos(s-a).
\end{align*}
Its maximum is attained exactly at
\(s\equiv a\pmod h\).  This is equivalent to \(z^d=-1\), again
giving \eqref{eq:equality-set}.  Multiplying the maximum of \(G_d\)
by two proves \eqref{eq:scalar}.
\end{proof}

\section{Proof of the upper bound}

\begin{proof}[Proof of Theorem~\ref{thm:main}]
Set
\[
  C_y=A_0+\omega^yA_1,
  \qquad
  U=A_0^\dagger A_1,
  \qquad
  M_y=\Id+\omega^yU.
\]
Then \(C_y=A_0M_y\).  The operator \(M_y\) is normal, since
\[
  M_y^\dagger M_y=M_yM_y^\dagger
  =
  2\Id+\omega^yU+\omega^{-y}U^\dagger.
\]
It follows that
\begin{equation}
  \absop{C_y}=\absop{M_y},
  \qquad
  \absop{C_y^\dagger}
  =
  A_0\absop{M_y}A_0^\dagger.
  \label{eq:moduli}
\end{equation}
Define
\[
  F_d(U)=\sum_{y=0}^{d-1}\absop{\Id+\omega^yU}.
\]
The continuous functional calculus for the unitary \(U\), together with
Lemma~\ref{lem:scalar}, yields
\begin{equation}
  0\leq F_d(U)\leq\lambda_d\Id.
  \label{eq:F-bound}
\end{equation}

Apply Lemma~\ref{lem:polar} to \(C_y\) and \(B_y\), sum over \(y\),
and use \eqref{eq:moduli}:
\begin{align*}
  \cI_d
  &\leq
  \frac12\sum_{y=0}^{d-1}
       \bigl(\absop{C_y}+\absop{C_y^\dagger}\bigr)\otimes\Id\\
  &=
  \frac12
  \bigl(F_d(U)+A_0F_d(U)A_0^\dagger\bigr)\otimes\Id\\
  &\leq \lambda_d\Id.
\end{align*}
This proves the operator upper bound.  The matching order-\(d\) strategy is
recalled and independently checked in Proposition~\ref{prop:lower}, completing
the proof of the equality.
\end{proof}

\subsection{An explicit positive-factor certificate}

The preceding proof can be packaged as a single exact identity.  Let
\(C_y=V_y\absop{C_y}\) be the polar decomposition, and set
\begin{align*}
  P_y
  &=
  \absop{C_y^\dagger}^{1/2}\otimes\Id
  -
  V_y\absop{C_y}^{1/2}\otimes B_y,\\
  G
  &=
  \bigl(\lambda_d\Id-F_d(U)\bigr)^{1/2}.
\end{align*}
Then
\begin{equation}
\boxed{
\begin{aligned}
  \lambda_d\Id-\cI_d
  &=
  \frac12\sum_{y=0}^{d-1}P_y^\dagger P_y\\
  &\quad+
  \frac12(G\otimes\Id)^\dagger(G\otimes\Id)\\
  &\quad+
  \frac12(GA_0^\dagger\otimes\Id)^\dagger
              (GA_0^\dagger\otimes\Id).
\end{aligned}
}
\label{eq:certificate}
\end{equation}
The square roots in \eqref{eq:certificate} are ordinary positive square roots
from functional calculus.  Thus \eqref{eq:certificate} is an exact
operator-level positive-factor certificate, not a finite-level NPA
approximation.

\section{Tightness and the order-\(d\) strategy}

The lower bound and the Weyl-based strategy are due to Perito et
al.~\cite{PeritoEtAl2026}.  For completeness, this section records a compact
polar form of the same saturating choice and proves directly that its Bob
operators satisfy the order relation.  This also makes the equality mechanism
in Lemma~\ref{lem:scalar} transparent.

\begin{proposition}[Self-contained lower certificate]
\label{prop:lower}
On \(\bbC^d\), let
\[
  Z\lvert j\rangle=\omega^j\lvert j\rangle,
  \qquad
  X\lvert j\rangle=\lvert j+1\bmod d\rangle,
\]
and use
\[
  A_0=Z,\qquad A_1=X,\qquad
  \lvert\Phi_d\rangle
  =
  \frac1{\sqrt d}\sum_{j=0}^{d-1}\lvert j,j\rangle.
\]
For \(0\leq y<d\), put
\begin{align}
  W_y&=\omega^yZ^\dagger X,\qquad
  H_y=\absop{\Id+W_y},\label{eq:WH}\\
  Q_y&=H_y^{-1}(\Id+W_y^\dagger)Z^\dagger,\qquad
  B_y=Q_y^T. \label{eq:Bob}
\end{align}
Then every \(B_y\) is unitary, \(B_y^d=\Id\), and
\[
  \langle\Phi_d\vert\cI_d\vert\Phi_d\rangle=\lambda_d.
\]
\end{proposition}

\begin{proof}
The Weyl relations give \(Z^d=X^d=\Id\) and \(ZX=\omega XZ\).
The action of \(W_y\) is a weighted cyclic shift:
\[
  W_y\lvert j\rangle
  =
  \omega^{\,y-(j+1\bmod d)}\lvert j+1\bmod d\rangle.
\]
Hence
\begin{equation}
  W_y^d=(-1)^{d-1}\Id,
  \qquad
  \spec(W_y)=
  \{\lambda:\lambda^d=(-1)^{d-1}\}.
  \label{eq:W-spectrum}
\end{equation}
Indeed, the first \(d\) orbit vectors of \(\lvert0\rangle\) are nonzero
multiples of the standard basis vectors.  Thus the minimal polynomial has
degree \(d\); because it divides
\(t^d-(-1)^{d-1}\), whose roots are distinct, the displayed spectrum is
simple and complete.
In particular, \(-1\notin\spec(W_y)\), so \(H_y\) is invertible.

Let \(C_y=Z+\omega^yX=Z(\Id+W_y)\).  Since \(W_y\) is normal,
\[
  C_y
  =
  \underbrace{Z(\Id+W_y)H_y^{-1}}_{V_y}\,H_y
\]
is its polar decomposition.  Thus \(Q_y=V_y^\dagger\), proving that
\(Q_y\), and hence \(B_y\), is unitary.

It remains to prove the order relation.  On the spectrum in
\eqref{eq:W-spectrum}, define
\[
  q(z)=\frac{1+\overline z}{\lvert1+z\rvert}.
\]
Then \(Q_y=q(W_y)Z^\dagger\).  Since
\(
Z^\dagger W_yZ=\omega^{-1}W_y
\),
\begin{equation}
  Q_y^d
  =
  \prod_{r=0}^{d-1}q(\omega^{-r}W_y).
  \label{eq:Q-product}
\end{equation}
For every eigenvalue \(\lambda\) of \(W_y\), the corresponding scalar in
\eqref{eq:Q-product} equals
\begin{align*}
  \prod_{r=0}^{d-1}q(\omega^{-r}\lambda)
  &=
  \frac{\prod_{r=0}^{d-1}(1+\omega^r\overline\lambda)}
       {\prod_{r=0}^{d-1}\lvert1+\omega^{-r}\lambda\rvert}\\
  &=
  \frac{1-(-\overline\lambda)^d}
       {\lvert1-(-\lambda)^d\rvert}
  =\frac2{2}=1,
\end{align*}
where we used
\(
\prod_{r=0}^{d-1}(1+a\omega^r)=1-(-a)^d
\)
and \(\lambda^d=(-1)^{d-1}\).  Therefore \(Q_y^d=\Id\) and
\(B_y^d=\Id\).  In fact, \(Q_y\) is a nonzero weighted \(d\)-cycle in
an eigenbasis of \(W_y\), so its characteristic polynomial is
\(t^d-1\); every \(d\)-th root occurs.

Finally,
\[
  \langle\Phi_d\vert M\otimes N\vert\Phi_d\rangle
  =
  \frac1d\Tr(MN^T).
\]
Since \(B_y^T=Q_y=V_y^\dagger\),
\[
  \langle\Phi_d\vert C_y\otimes B_y\vert\Phi_d\rangle
  =
  \frac1d\Tr(C_yV_y^\dagger)
  =
  \frac1d\Tr(H_y),
\]
which is real and nonnegative.  Multiplication by \(\omega^y\) permutes
the spectrum in \eqref{eq:W-spectrum}, so every \(H_y\) has the same trace.
For any \(\lambda_0^d=(-1)^{d-1}\),
\[
  \Tr(H_0)
  =
  \sum_{r=0}^{d-1}\lvert1+\omega^r\lambda_0\rvert
  =
  \lambda_d
\]
by the equality case of Lemma~\ref{lem:scalar}.  Summing over \(y\)
therefore gives
\(
d^{-1}\sum_y\Tr(H_y)=\lambda_d
\),
as claimed.
\end{proof}

\begin{remark}
Because \(C_y\) is invertible in this construction, equality in the trace
polar bound fixes \(B_y^T\) to be \(V_y^\dagger\).  The Bob observables in
\eqref{eq:Bob} therefore give a polar-decomposition representation of an
attaining strategy.  The strategy in Ref.~\cite{PeritoEtAl2026} reaches the
sum of these termwise trace bounds, so the same equality condition identifies
its Bob observables with this polar form.  No new lower-bound strategy is
claimed here.
\end{remark}

\section{Consequences and scope}

\begin{corollary}[The barred functional]
\label{cor:barred}
With the convention
\[
  \overline{\cI}_d
  =
  \cI_d+\Reop(A_0\otimes B_d),
\]
where \(B_d\) is a further unitary, its quantum value is
\[
  \beta_{\mathrm Q}(\overline{\cI}_d)
  =
  2\csc\!\left(\frac{\pi}{2d}\right)+1.
\]
\end{corollary}

\begin{proof}
Theorem~\ref{thm:main} and
\(
\Reop(A_0\otimes B_d)\leq\Id
\)
give the upper bound.  Proposition~\ref{prop:lower}, augmented by
\(
B_d=Z^\dagger
\),
saturates the extra term on \(\lvert\Phi_d\rangle\).
\end{proof}

This corollary determines a Bell value.  It does not prove uniqueness of the
maximizing correlation, self-testing, or the maximal-global-randomness claim
posed separately in Ref.~\cite{PeritoEtAl2026}.  Those questions require
additional rigidity information not supplied by the value alone.

For the first five dimensions, Theorem~\ref{thm:main} gives
\[
\begin{array}{c@{\qquad}c}
\toprule
d&\beta_{\mathrm Q}(\cI_d)\\
\midrule
2&2\sqrt2\\
3&4\\
4&2\sqrt{4+2\sqrt2}\\
5&2(1+\sqrt5)\\
6&2(\sqrt6+\sqrt2)\\
\bottomrule
\end{array}
\]
The \(d=2\) case is the usual CHSH value; the theorem supplies one uniform
argument for all \(d\).

\section{Verification artifacts}

The accompanying source package contains:
\begin{itemize}[leftmargin=1.5em]
  \item a deterministic verifier for the finite-dimensional matrix identities,
        including deliberately singular polar factors;
  \item exact symbolic checks of the displayed low-dimensional radical values
        and finite trigonometric identities;
  \item matrix checks of the Weyl spectrum, Bob unitarity and order,
        saturation, and the operator certificate over a documented range of
        dimensions;
  \item a machine-readable claim and dependency ledger; and
  \item SHA-256 hashes for the released sources and rendered paper.
\end{itemize}
These checks are secondary validation of the analytic proof.  They are not a
formal proof assistant verification of the all-dimensional theorem, and the
floating-point matrix tests cannot substitute for the arguments above.

\section*{Priority, attribution, and disclosure}

The Bell family, the conjectured formula, the Weyl construction, the
admissible Bob observables, and the matching lower bound are due to Perito,
D'Avino, Jung, Mironowicz, Acin, and Augusiak
\cite{PeritoEtAl2026}.  The narrow contribution claimed here is the analytic
upper bound, including its arbitrary-unitary strengthening and the associated
operator certificate.

As of 26 July 2026, a focused search found no public proof of Conjecture~1 and
no older theorem located in that search that directly supplies this bound.
Because the conjecture is recent, that search cannot rule out unindexed,
unpublished, or concurrent work.  Any priority statement is therefore made
only ``to the best of our knowledge.''

The proof architecture, derivations, manuscript, verification code, and
literature audit were developed with heavy assistance from ChatGPT 5.6 Sol
under Alec Kriebel's direction.  No originating author or other researcher was
contacted.  The work is public to enable checking and remains unrefereed.

\begin{thebibliography}{9}

\bibitem{PeritoEtAl2026}
I. Perito, R. D'Avino, M. Jung, P. Mironowicz, A. Ac\'in, and R. Augusiak,
\newblock Bell inequalities tailored to optimal global randomness
  certification,
\newblock \emph{arXiv:2606.21362v3} (2026),
\newblock \url{https://arxiv.org/abs/2606.21362}.

\bibitem{Tsirelson1980}
B. S. Tsirelson,
\newblock Quantum generalizations of Bell's inequality,
\newblock \emph{Letters in Mathematical Physics} \textbf{4}, 93--100 (1980).

\bibitem{NPA2007}
M. Navascu\'es, S. Pironio, and A. Ac\'in,
\newblock Bounding the set of quantum correlations,
\newblock \emph{Physical Review Letters} \textbf{98}, 010401 (2007).

\bibitem{NPA2008}
M. Navascu\'es, S. Pironio, and A. Ac\'in,
\newblock A convergent hierarchy of semidefinite programs characterizing the
  set of quantum correlations,
\newblock \emph{New Journal of Physics} \textbf{10}, 073013 (2008).

\end{thebibliography}

\end{document}
